% generated by GAPDoc2LaTeX from XML source (Frank Luebeck)
\documentclass[a4paper,11pt]{report}

\usepackage[top=37mm,bottom=37mm,left=27mm,right=27mm]{geometry}
\sloppy
\pagestyle{myheadings}
\usepackage{amssymb}
\usepackage[utf8]{inputenc}
\usepackage{makeidx}
\makeindex
\usepackage{color}
\definecolor{FireBrick}{rgb}{0.5812,0.0074,0.0083}
\definecolor{RoyalBlue}{rgb}{0.0236,0.0894,0.6179}
\definecolor{RoyalGreen}{rgb}{0.0236,0.6179,0.0894}
\definecolor{RoyalRed}{rgb}{0.6179,0.0236,0.0894}
\definecolor{LightBlue}{rgb}{0.8544,0.9511,1.0000}
\definecolor{Black}{rgb}{0.0,0.0,0.0}

\definecolor{linkColor}{rgb}{0.0,0.0,0.554}
\definecolor{citeColor}{rgb}{0.0,0.0,0.554}
\definecolor{fileColor}{rgb}{0.0,0.0,0.554}
\definecolor{urlColor}{rgb}{0.0,0.0,0.554}
\definecolor{promptColor}{rgb}{0.0,0.0,0.589}
\definecolor{brkpromptColor}{rgb}{0.589,0.0,0.0}
\definecolor{gapinputColor}{rgb}{0.589,0.0,0.0}
\definecolor{gapoutputColor}{rgb}{0.0,0.0,0.0}

%%  for a long time these were red and blue by default,
%%  now black, but keep variables to overwrite
\definecolor{FuncColor}{rgb}{0.0,0.0,0.0}
%% strange name because of pdflatex bug:
\definecolor{Chapter }{rgb}{0.0,0.0,0.0}
\definecolor{DarkOlive}{rgb}{0.1047,0.2412,0.0064}


\usepackage{fancyvrb}

\usepackage{mathptmx,helvet}
\usepackage[T1]{fontenc}
\usepackage{textcomp}


\usepackage[
            pdftex=true,
            bookmarks=true,        
            a4paper=true,
            pdftitle={Written with GAPDoc},
            pdfcreator={LaTeX with hyperref package / GAPDoc},
            colorlinks=true,
            backref=page,
            breaklinks=true,
            linkcolor=linkColor,
            citecolor=citeColor,
            filecolor=fileColor,
            urlcolor=urlColor,
            pdfpagemode={UseNone}, 
           ]{hyperref}

\newcommand{\maintitlesize}{\fontsize{50}{55}\selectfont}

% write page numbers to a .pnr log file for online help
\newwrite\pagenrlog
\immediate\openout\pagenrlog =\jobname.pnr
\immediate\write\pagenrlog{PAGENRS := [}
\newcommand{\logpage}[1]{\protect\write\pagenrlog{#1, \thepage,}}
%% were never documented, give conflicts with some additional packages

\newcommand{\GAP}{\textsf{GAP}}

%% nicer description environments, allows long labels
\usepackage{enumitem}
\setdescription{style=nextline}

%% depth of toc
\setcounter{tocdepth}{1}





%% command for ColorPrompt style examples
\newcommand{\gapprompt}[1]{\color{promptColor}{\bfseries #1}}
\newcommand{\gapbrkprompt}[1]{\color{brkpromptColor}{\bfseries #1}}
\newcommand{\gapinput}[1]{\color{gapinputColor}{#1}}


\begin{document}

\logpage{[ 0, 0, 0 ]}
\begin{titlepage}
\mbox{}\vfill

\begin{center}{\maintitlesize \textbf{ liebreadth \\
\mbox{}}}\\
\vfill

\hypersetup{pdftitle={ liebreadth }}
\markright{\scriptsize \mbox{}\hfill  liebreadth  \hfill\mbox{}}
{\Huge \textbf{ Breadth in finite dimensional algebras \\
\mbox{}}}\\
\vfill

{\Huge  1.0 \mbox{}}\\[1cm]
{ 24 September 2026 \mbox{}}\\[1cm]
\mbox{}\\[2cm]
{\Large \textbf{\strut  {\a'O}scar Fern{\a'a}ndez Ayala    \strut\mbox{}}}\\
\hypersetup{pdfauthor={ {\a'O}scar Fern{\a'a}ndez Ayala   }}
\end{center}\vfill

\mbox{}\\
{\mbox{}\\
\small \noindent \textbf{ {\a'O}scar Fern{\a'a}ndez Ayala   }  Email: \href{mailto://oscar00ayala@gmail.com} {\texttt{oscar00ayala@gmail.com}}\\
  Homepage: \href{https://Osferay.github.io} {\texttt{https://Osferay.github.io}}}\\
\end{titlepage}

\newpage\setcounter{page}{2}
\newpage

\def\contentsname{Contents\logpage{[ 0, 0, 1 ]}}

\tableofcontents
\newpage

     
\chapter{\textcolor{Chapter }{Preface}}\label{Chapter_preface}
\logpage{[ 1, 0, 0 ]}
\hyperdef{L}{X874E1D45845007FE}{}
{
  In this package we compute the breadth of Lie algebras. The details of the
algorithms are available in \cite{breadth}. This package also includes functions to compute covered Lie algebras of
maximal class using a process called inflation, which is described by Caranti,
Mattarei, and Newman \cite{inflation}. 
\section{\textcolor{Chapter }{Introduction}}\label{Chapter_Preface_Section_Introduction}
\logpage{[ 1, 1, 0 ]}
\hyperdef{L}{X7DFB63A97E67C0A1}{}
{
  Given a finite\texttt{\symbol{45}}dimensional algebra $L$, we define its lower central series as $L = L^1 > L^2 > \dots$, where $L^{i+1} = L^i L$. The algebra $L$ is nilpotent if there exists $c \in \mathbb{N}$ such that $L^{c+1} = \{0\}$, and the minimal such $c$ is called the class $\mathrm{cl}(L)$ of $L$. For a nilpotent algebra $L$, the type of $L$ is the vector $(d_1,\dots,d_c)$, where $d_i = \dim L^i$. A nilpotent algebra $L$ is said to be of maximal class if the type of $L$ is $(2,1,\dots,1)$. The centralizer of $x \in L$ is the subspace of $L$ defined by $C_L(x) = \{a \in L \mid ax = 0\}$. For an algebra $L$, we define 
\[ \mathrm{br}(L) = \max \{ \mathrm{br}(x) \mid x \in L\}, \;\; \mbox{ where }
\;\; \mathrm{br}(x) = \dim(L) - \dim(C_L(x)). \]
 The class\texttt{\symbol{45}}breadth conjecture asserts that $\mathrm{cl}(L) \leq \mathrm{br}(L)+1$ for an algebra $L$. The conjecture holds for nilpotent Lie algebras over infinite fields and for
nilpotent associative algebras over arbitrary fields. 

 Let $L$ be a Lie algebra over a field $K$, and let $B = \{ b_1, \ldots, b_n \}$ be a basis of $L$. The multiplication in $L$ is described by structure constants defined by 
\[ b_i \cdot b_j = \sum_{k=1}^n c_{ijk} b_k \quad\mbox{ for } 1\leq i,j \leq n. \]
 For each $k$, write $C_k = (c_{ijk})_{1 \leq i,j \leq n}$ for the $n \times n$ matrix over $K$ with entries $c_{ijk}$, and these matrices are the structure matrices of $L$. For $x = x_1 b_1 + \ldots + x_n b_n \in L$, let $\overline{x} = (x_1, \ldots, x_n) \in K^n$ denote the coefficient vector of $x$. Then $C_k \overline{x}^{tr}$ is a column vector over $K$. We write $M_B(x)$ for the $n \times n$ matrix over $K$ whose $k$\texttt{\symbol{45}}th column is $C_k \overline{x}^{\mathrm{tr}}$, which is the adjoint matrix of $L$ with respect to the basis $B$. }

 }

   
\chapter{\textcolor{Chapter }{Lie algebras and breadth}}\label{Chapter_lie}
\logpage{[ 2, 0, 0 ]}
\hyperdef{L}{X8180FB547D5676E4}{}
{
  
\section{\textcolor{Chapter }{Attributes}}\label{Chapter_Lie_algebras_and_breadth_Section_Attributes}
\logpage{[ 2, 1, 0 ]}
\hyperdef{L}{X7C701DBF7BAE649A}{}
{
  

\subsection{\textcolor{Chapter }{LieClass (for IsLieNilpotent)}}
\logpage{[ 2, 1, 1 ]}\nobreak
\hyperdef{L}{X829CF8CB7B08852D}{}
{\noindent\textcolor{FuncColor}{$\triangleright$\enspace\texttt{LieClass({\mdseries\slshape L})\index{LieClass@\texttt{LieClass}!for IsLieNilpotent}
\label{LieClass:for IsLieNilpotent}
}\hfill{\scriptsize (attribute)}}\\


 Computes the class of the nilpotent Lie algebra $L$. }

 

\subsection{\textcolor{Chapter }{LieType (for IsLieNilpotent)}}
\logpage{[ 2, 1, 2 ]}\nobreak
\hyperdef{L}{X86E34A087C73F4D4}{}
{\noindent\textcolor{FuncColor}{$\triangleright$\enspace\texttt{LieType({\mdseries\slshape L})\index{LieType@\texttt{LieType}!for IsLieNilpotent}
\label{LieType:for IsLieNilpotent}
}\hfill{\scriptsize (attribute)}}\\


 Computes the type of the nilpotent Lie algebra $L$. }

 

\subsection{\textcolor{Chapter }{IsOfMaximalClass (for IsLieNilpotent)}}
\logpage{[ 2, 1, 3 ]}\nobreak
\hyperdef{L}{X7CC91C177F001E7C}{}
{\noindent\textcolor{FuncColor}{$\triangleright$\enspace\texttt{IsOfMaximalClass({\mdseries\slshape L})\index{IsOfMaximalClass@\texttt{IsOfMaximalClass}!for IsLieNilpotent}
\label{IsOfMaximalClass:for IsLieNilpotent}
}\hfill{\scriptsize (property)}}\\
\textbf{\indent Returns:\ }
\texttt{true} or \texttt{false} 



 Returns whether the nilpotent Lie algebra $L$ is of maximal class or not. }

 

\subsection{\textcolor{Chapter }{StructureMatrices (for IsLieAlgebra)}}
\logpage{[ 2, 1, 4 ]}\nobreak
\hyperdef{L}{X855BD6DB7CA138C6}{}
{\noindent\textcolor{FuncColor}{$\triangleright$\enspace\texttt{StructureMatrices({\mdseries\slshape L})\index{StructureMatrices@\texttt{StructureMatrices}!for IsLieAlgebra}
\label{StructureMatrices:for IsLieAlgebra}
}\hfill{\scriptsize (attribute)}}\\


 Computes the structure matrices of the nilpotent Lie algebra $L$. }

 

\subsection{\textcolor{Chapter }{BasisLieCenter (for IsLieAlgebra)}}
\logpage{[ 2, 1, 5 ]}\nobreak
\hyperdef{L}{X84C91702821B9D8F}{}
{\noindent\textcolor{FuncColor}{$\triangleright$\enspace\texttt{BasisLieCenter({\mdseries\slshape L})\index{BasisLieCenter@\texttt{BasisLieCenter}!for IsLieAlgebra}
\label{BasisLieCenter:for IsLieAlgebra}
}\hfill{\scriptsize (attribute)}}\\


 Returns the basis elements of the Lie algebra $L$ that are contained in the center of $L$. }

 

\subsection{\textcolor{Chapter }{BasisLieDerived (for IsLieAlgebra)}}
\logpage{[ 2, 1, 6 ]}\nobreak
\hyperdef{L}{X79539FB282963F71}{}
{\noindent\textcolor{FuncColor}{$\triangleright$\enspace\texttt{BasisLieDerived({\mdseries\slshape L})\index{BasisLieDerived@\texttt{BasisLieDerived}!for IsLieAlgebra}
\label{BasisLieDerived:for IsLieAlgebra}
}\hfill{\scriptsize (attribute)}}\\


 Returns the basis elements of the Lie algebra $L$ that are contained in the derived subalgebra of $L$. }

 

\subsection{\textcolor{Chapter }{LieAdjointMatrix (for IsLieAlgebra)}}
\logpage{[ 2, 1, 7 ]}\nobreak
\hyperdef{L}{X794FAF3F8277235F}{}
{\noindent\textcolor{FuncColor}{$\triangleright$\enspace\texttt{LieAdjointMatrix({\mdseries\slshape L})\index{LieAdjointMatrix@\texttt{LieAdjointMatrix}!for IsLieAlgebra}
\label{LieAdjointMatrix:for IsLieAlgebra}
}\hfill{\scriptsize (attribute)}}\\


 Computes the adjoint matrix of the Lie algebra $L$. }

 

\subsection{\textcolor{Chapter }{PrintLiePresentation (for IsLieAlgebra)}}
\logpage{[ 2, 1, 8 ]}\nobreak
\hyperdef{L}{X80F94DFF78B2777A}{}
{\noindent\textcolor{FuncColor}{$\triangleright$\enspace\texttt{PrintLiePresentation({\mdseries\slshape L})\index{PrintLiePresentation@\texttt{PrintLiePresentation}!for IsLieAlgebra}
\label{PrintLiePresentation:for IsLieAlgebra}
}\hfill{\scriptsize (attribute)}}\\


 Prints the Lie presentation of the Lie algebra $L$. }

 }

 
\section{\textcolor{Chapter }{Breadth}}\label{Chapter_Lie_algebras_and_breadth_Section_Breadth}
\logpage{[ 2, 2, 0 ]}
\hyperdef{L}{X8648E5BC79C8AEA8}{}
{
  

\subsection{\textcolor{Chapter }{InfoLieBreadth}}
\logpage{[ 2, 2, 1 ]}\nobreak
\hyperdef{L}{X82F4EEAD7DFC6470}{}
{\noindent\textcolor{FuncColor}{$\triangleright$\enspace\texttt{InfoLieBreadth\index{InfoLieBreadth@\texttt{InfoLieBreadth}}
\label{InfoLieBreadth}
}\hfill{\scriptsize (info class)}}\\


 Info class for the functions of the breadth of Lie algebras. }

 

\subsection{\textcolor{Chapter }{LieBreadth (for IsLieNilpotent)}}
\logpage{[ 2, 2, 2 ]}\nobreak
\hyperdef{L}{X7C093B047BA37D4A}{}
{\noindent\textcolor{FuncColor}{$\triangleright$\enspace\texttt{LieBreadth({\mdseries\slshape L})\index{LieBreadth@\texttt{LieBreadth}!for IsLieNilpotent}
\label{LieBreadth:for IsLieNilpotent}
}\hfill{\scriptsize (attribute)}}\\


 Computes the breadth of the Lie algebra $L$. }

 

\subsection{\textcolor{Chapter }{IsTrueClassBreadth (for IsLieNilpotent)}}
\logpage{[ 2, 2, 3 ]}\nobreak
\hyperdef{L}{X86CB363980B40569}{}
{\noindent\textcolor{FuncColor}{$\triangleright$\enspace\texttt{IsTrueClassBreadth({\mdseries\slshape L})\index{IsTrueClassBreadth@\texttt{IsTrueClassBreadth}!for IsLieNilpotent}
\label{IsTrueClassBreadth:for IsLieNilpotent}
}\hfill{\scriptsize (property)}}\\
\textbf{\indent Returns:\ }
\texttt{true} or \texttt{false} 



 Returns whether the class\texttt{\symbol{45}}breadth conjecture holds for the
Lie algebra $L$. }

 

\subsection{\textcolor{Chapter }{TGroupBreadth}}
\logpage{[ 2, 2, 4 ]}\nobreak
\hyperdef{L}{X7F177A627B2917E9}{}
{\noindent\textcolor{FuncColor}{$\triangleright$\enspace\texttt{TGroupBreadth({\mdseries\slshape G})\index{TGroupBreadth@\texttt{TGroupBreadth}}
\label{TGroupBreadth}
}\hfill{\scriptsize (function)}}\\


 Computes the breadth of the $\mathcal{T}$\texttt{\symbol{45}}group $G$. }

 }

 
\section{\textcolor{Chapter }{Inflation of Lie algebras}}\label{Chapter_Lie_algebras_and_breadth_Section_Inflation_of_Lie_algebras}
\logpage{[ 2, 3, 0 ]}
\hyperdef{L}{X79E2250284B5313A}{}
{
  A grading for a Lie algebra $L$ is a decomposition $L = \bigoplus_{i=1}^n L_i$ that respects the Lie bracket, \emph{i.e.} $[L_i,L_j]\subseteq L_{i+j}$. All nilpotent Lie algebras can be graded by taking $L_i = \gamma_i(L)/\gamma_{i+1}(L)$. Let $L$ be a nilpotent Lie algebra of maximal class. The two\texttt{\symbol{45}}step
centralizers are the sets $C_i = C_{L_1}(L_i) = \{x \in L_1 \mid [x,L_i] = 0 \}$ for all $2\leq i \leq c$. Let $\mathcal{C}=\{C_i\}\setminus L_1$. We say that a Lie algebra of maximal class is covered if the set $\mathcal{C}$ consist of all one\texttt{\symbol{45}}dimensional subspaces of $L_1$. 

 Let $L = \bigoplus_{i=1}^n L_i$ be a graded Lie algebra over $K=\mathbb{F}_q$ for some prime power $q=p^n$. Consider the field extension $A=K[\varepsilon]/\langle \varepsilon^p \rangle$ which is a vector space over $K$ of dimension $p$. The algebra $A$ is an associative, commutative algebra with a unit. The algebra $L\otimes A$ over $K$, defined by $[x\otimes a, y\otimes b] = [x,y]\otimes ab$, is a graded Lie algebra. Let $M$ be a maximal ideal of $L$ and consider the Lie subalgebra $M^\uparrow = M\otimes A$. Let $D$ be a derivation of $L$ of degree 1. Define $D^\uparrow$ by: $D^\uparrow(x\otimes \varepsilon^i)=D(x)\otimes \varepsilon^i$, which is a derivation of $M^\uparrow$ of degree 1. Define the derivation $E\in \mathrm{Der} M^\uparrow$ by $E(x\otimes \varepsilon^i) =
D^\uparrow(x\otimes\varepsilon^i\cdot\varepsilon^{p-1})+1\otimes\partial_\varepsilon(\varepsilon^i)$. For $s\in L_1\setminus M$, take $D=\mathrm{ad}_s$, and extend it naturally to $M^\uparrow$. Denote $E_{s^\prime}$ the derivation defined above. The \emph{inflation} ${}^{M}L$ of $L$ at $M$ by $s\in L_1\setminus M$ is the graded Lie algebra obtained as an extension of $M^\uparrow$ by an element $s^\prime$, which is the extension of $s$ that induces the derivation $E_{s^\prime}$, that is, 
\[ [x\otimes a, s^\prime] = E_{s^\prime}(x\otimes a) =
\mathrm{ad}_{s^\prime}(x\otimes\varepsilon^i)+x\otimes\partial_\varepsilon(\varepsilon^i). \]
 

\subsection{\textcolor{Chapter }{InfoInflation}}
\logpage{[ 2, 3, 1 ]}\nobreak
\hyperdef{L}{X7B546C44790587C4}{}
{\noindent\textcolor{FuncColor}{$\triangleright$\enspace\texttt{InfoInflation\index{InfoInflation@\texttt{InfoInflation}}
\label{InfoInflation}
}\hfill{\scriptsize (info class)}}\\


 Info class for the functions of the inflation of Lie algebras. }

 

\subsection{\textcolor{Chapter }{LieNilpotentGrading (for IsLieNilpotent)}}
\logpage{[ 2, 3, 2 ]}\nobreak
\hyperdef{L}{X821797467A3A523E}{}
{\noindent\textcolor{FuncColor}{$\triangleright$\enspace\texttt{LieNilpotentGrading({\mdseries\slshape L})\index{LieNilpotentGrading@\texttt{LieNilpotentGrading}!for IsLieNilpotent}
\label{LieNilpotentGrading:for IsLieNilpotent}
}\hfill{\scriptsize (attribute)}}\\


 Computes a grading of the nilpotent Lie algebra $L$ using the lower central series. }

 

\subsection{\textcolor{Chapter }{LieTwoStepCentralizers (for IsLieNilpotent)}}
\logpage{[ 2, 3, 3 ]}\nobreak
\hyperdef{L}{X85CE14C17816A97D}{}
{\noindent\textcolor{FuncColor}{$\triangleright$\enspace\texttt{LieTwoStepCentralizers({\mdseries\slshape L})\index{LieTwoStepCentralizers@\texttt{LieTwoStepCentralizers}!for IsLieNilpotent}
\label{LieTwoStepCentralizers:for IsLieNilpotent}
}\hfill{\scriptsize (attribute)}}\\


 Computes the two\texttt{\symbol{45}}step centralizers of the Lie algebra $L$. }

 

\subsection{\textcolor{Chapter }{IsLieCovered (for IsLieNilpotent)}}
\logpage{[ 2, 3, 4 ]}\nobreak
\hyperdef{L}{X7E24E6DF824913CD}{}
{\noindent\textcolor{FuncColor}{$\triangleright$\enspace\texttt{IsLieCovered({\mdseries\slshape L})\index{IsLieCovered@\texttt{IsLieCovered}!for IsLieNilpotent}
\label{IsLieCovered:for IsLieNilpotent}
}\hfill{\scriptsize (property)}}\\
\textbf{\indent Returns:\ }
\texttt{true} or \texttt{false} 



 Returns whether the Lie algebra $L$ is covered or not. }

 

\subsection{\textcolor{Chapter }{PolynomialAlgebra (for IsField and IsFinite)}}
\logpage{[ 2, 3, 5 ]}\nobreak
\hyperdef{L}{X7F13AC8581C578E0}{}
{\noindent\textcolor{FuncColor}{$\triangleright$\enspace\texttt{PolynomialAlgebra({\mdseries\slshape F})\index{PolynomialAlgebra@\texttt{PolynomialAlgebra}!for IsField and IsFinite}
\label{PolynomialAlgebra:for IsField and IsFinite}
}\hfill{\scriptsize (attribute)}}\\


 For a finite field $F$, computes the polynomial algebra $F[x]$. }

 

\subsection{\textcolor{Chapter }{LieCoveredInflated (for IsInt)}}
\logpage{[ 2, 3, 6 ]}\nobreak
\hyperdef{L}{X86B86B1784C39326}{}
{\noindent\textcolor{FuncColor}{$\triangleright$\enspace\texttt{LieCoveredInflated({\mdseries\slshape n})\index{LieCoveredInflated@\texttt{LieCoveredInflated}!for IsInt}
\label{LieCoveredInflated:for IsInt}
}\hfill{\scriptsize (attribute)}}\\


 For $n=2,3,4$, computes the covered Lie algebras of maximal class using the polynomial
algebra of dimension $n$. For $n=3$, one can directly load the Lie algebra by reading the file
\texttt{\symbol{92}}texttt\texttt{\symbol{123}}inflation\texttt{\symbol{92}}{\textunderscore}3.g\texttt{\symbol{125}}. }

 

\subsection{\textcolor{Chapter }{LieMinimalQuotientClassBreadth (for IsLieAlgebra)}}
\logpage{[ 2, 3, 7 ]}\nobreak
\hyperdef{L}{X828D60C387E6060B}{}
{\noindent\textcolor{FuncColor}{$\triangleright$\enspace\texttt{LieMinimalQuotientClassBreadth({\mdseries\slshape L})\index{LieMinimalQuotientClassBreadth@\texttt{LieMinimalQuotientClassBreadth}!for IsLieAlgebra}
\label{LieMinimalQuotientClassBreadth:for IsLieAlgebra}
}\hfill{\scriptsize (attribute)}}\\


 Given a covered Lie algebras of maximal class $L$, computes the minimal quotient for which the class\texttt{\symbol{45}}breadth
conjecture does not hold. }

 }

 }

 \def\bibname{References\logpage{[ "Bib", 0, 0 ]}
\hyperdef{L}{X7A6F98FD85F02BFE}{}
}

\bibliographystyle{alpha}
\bibliography{bib.xml}

\addcontentsline{toc}{chapter}{References}

\def\indexname{Index\logpage{[ "Ind", 0, 0 ]}
\hyperdef{L}{X83A0356F839C696F}{}
}

\cleardoublepage
\phantomsection
\addcontentsline{toc}{chapter}{Index}


\printindex

\immediate\write\pagenrlog{["Ind", 0, 0], \arabic{page},}
\immediate\write\pagenrlog{["Ind", 0, 0], \arabic{page},}
\immediate\write\pagenrlog{["Ind", 0, 0], \arabic{page},}
\immediate\write\pagenrlog{["Ind", 0, 0], \arabic{page},}
\newpage
\immediate\write\pagenrlog{["End"], \arabic{page}];}
\immediate\closeout\pagenrlog
\end{document}
